\begin{problem}[Warmup]{(a)}
    Give examples of the following group homomorphisms:
    \begin{enumerate}
        \item[(a)] A one-to-one homomorphism that is not an isomorphism.
        \item[(b)] An onto homomorphism that is not an isomorphism.
        \item[(c)] A homomorphism that is neither one-to-one nor onto.
    \end{enumerate}
\end{problem}
\unsolved

\begin{problem}[Warmup]{(b)}
    Suppose that $\varphi$ is a group homomorphism with domain $\mathbb{Z}_8$. What are the possible orders of $\varphi(1)$? What about $\varphi(4)$?
\end{problem}
\unsolved

\begin{problem}[Warmup]{(c)}
    Suppose that $g$ has infinite order and $\varphi$ is a group homomorphism. Need $\varphi(g)$ have infinite order? What if $\varphi$ is an isomorphism?
\end{problem}
\unsolved

\begin{problem}[Warmup]{(d)}
    Explain why none of the following groups are isomorphic (even though all have eight elements): \[\mathbb{Z}_8, \quad \mathbb{Z}_4 \times \mathbb{Z}_2, \quad D_4.\]
\end{problem}
\unsolved

\begin{problem}[Warmup]{(e)}
    Explain why $\mathbb{Z}[i]$ and $\mathbb{Z}$ are not isomorphic as additive groups. That is, why isn't $\mathbb{Z}[i]$ cyclic?
\end{problem}
\unsolved

\begin{problem}[Warmup]{(f)}
    Can a group be isomorphic to a proper subgroup of itself? What if the group is finite?
\end{problem}
\unsolved

\begin{problem}[Warmup]{(g)}
    Let $R$ be a finite ring, and consider its additive group and its group of units. Could these two groups be isomorphic?
\end{problem}
\unsolved

\begin{problem}{1}
    Consider the set of complex numbers $\{1, -1, i, -i\}$. Note this is a group under
    multiplication. Show that this group is isomorphic to $\mathbb{Z}_4$.
\end{problem}
\unsolved

\begin{problem}{2}
    Use the Internal Characterization Theorem 28.3 to show that $\mathbb{Z}_{341}$ is
    isomorphic to $\mathbb{Z}_{11} \times \mathbb{Z}_{31}$. Then specify an explicit
    isomorphism.
\end{problem}
\unsolved

\begin{problem}{3}
    Suppose that $G, H$, and $K$ are groups. Prove that the direct products
    \[(G \times H) \times K \quad\text{and}\quad G \times (H \times K)\]
    are isomorphic. For this reason, we usually omit the parentheses when describing such
    groups.
\end{problem}
\unsolved

\begin{problem}{4}
    Complete the proof the Internal Characterization Theorem 28.3; that is, show that the
    function $\psi$ defined from $G_1 \times G_2$ to $G$ is in fact an isomorphism.
\end{problem}
\unsolved

\begin{problem}{5}
    If $G$ is a group and $g$ is some fixed element of $G$, show that the map $\varphi_g$
    defined by $\varphi_g(x) = gxg^{-1}$, for all $x \in G$, is an isomorphism from $G$
    onto itself.
\end{problem}
\unsolved

\begin{problem}{6}
    Here you'll prove parts a and b of Theorem 28.1. Let $\varphi : G \to H$ be a group
    homomorphism between the groups $G$ and $H$, where $1_G$ and $1_H$ are their respective
    identities.
    \begin{enumerate}
        \item[(a)] Prove that $\varphi(1_G) = 1_H$.
        \item[(b)] For $g \in G$, prove that ${(\varphi(g))}^{-1} = \varphi(g^{-1})$.
    \end{enumerate}
\end{problem}
\unsolved

\begin{problem}{7}
    Explain why $\mathbb{Q}$ and $\mathbb{Z}$ are not isomorphic as additive groups. That
    is, why isn't $\mathbb{Q}$ a cyclic group?
\end{problem}
\unsolved

\begin{problem}{8}
    For group $G$, consider the map $\varphi : G \to G$ given by $\varphi(g) = g^{-1}$.
    Show $\varphi$ is an isomorphism if and only if $G$ is abelian.
\end{problem}
\unsolved

\begin{problem}{9}
    Let $a, b, c \in \mathbb{R}$, not all zero.
    \begin{enumerate}
        \item[(a)] Show that
        \[P = \{(x, y, z) \in \mathbb{R}^3 : ax + by + cz = 0\}\]
        is a subgroup of $\mathbb{R}^3$.
        \item[(b)] Show that
        \[L = \{(ak, bk, ck) \in \mathbb{R}^3 : k \in \mathbb{R}\}\]
        is a subgroup of $\mathbb{R}^3$.
        \item[(c)] Use the Internal Characterization Theorem to show that $\mathbb{R}^3$ is
        isomorphic to $P \times L$. What does this mean geometrically?
    \end{enumerate}
\end{problem}
\unsolved

\begin{problem}{10}
    By way of analogy with Example 28.4, show that the additive group
    $\mathbb{Z}[\sqrt{2}]$ is isomorphic to a direct product of two non-trivial groups.
\end{problem}
\unsolved

\begin{problem}{11}
    Use the Internal Characterization Theorem to show that $U(\mathbb{Z}_{15})$ is
    isomorphic to a direct product of two non-trivial groups.
\end{problem}
\unsolved

\begin{problem}{12}
    Let $n$ be a positive integer, and consider the group $G_n$ described in
    Exercise 24.15. We will relabel these elements in such a way that we can consider
    $G_n$ as the list of objects
    \[G_n = \{I, R, R^2, \ldots, R^{n - 1}, F, FR, FR^2, \ldots, FR^{n - 1}\},\]
    where $I$ is the label we assign to the identity, and the elements $(1, [k])$ are
    labeled as $R^k$, and the elements $(-1, [k])$ are labeled as $FR^k$.
    \begin{enumerate}
        \item[(a)] Show that the elements of $G_n$ under this labeling do satisfy the
        identities
        \[RF = FR^{n - 1}, \quad F^2 = 1, \quad\text{and}\quad R^n = 1.\]
        In fact, it is possible to prove that our list above of the elements $R^k, FR^k$,
        together with these identities, abstractly characterizes the group $G_n$. We will
        not rigorously prove this here, however.
        \item[(b)] Using the abstract characterization of $G_n$ described above, prove that
        the dihedral group $D_n$ as discussed in Chapter 22, is isomorphic to $G_n$. (You
        need only establish a one-to-one correspondence that preserves the identities.)
        Henceforth, when we refer to the dihedral groups $D_n$, we will often use the
        notation of this exercise.
    \end{enumerate}
\end{problem}
\unsolved

\begin{problem}{13}
    In the group $D_3$, let $H_1 = \{I, F\}$ and $H_2 = \{I, R, R^2\}$. (Note that we are
    using the notation of the previous exercise.) Determine which of the four criteria of
    the Internal Characterization Theorem are satisfied by these two subgroups.
\end{problem}
\unsolved

\begin{problem}{14}
    In this exercise we introduce a more efficient and compact notation for the elements
    of the group $Q_8$ of quaternions, discussed as Example 24.15. We will relabel these
    elements in such a way that we can consider $Q_8$ as the list of objects
    \[\{1, -1, i, -i, j, -j, k, -k\},\]
    where 1 is the label we assign to the identity matrix, $i$ is the label assigned to
    matrix $\begin{pmatrix} 0 & 1 \\ -1 & 0 \end{pmatrix}$, $j$ is the label assigned to
    matrix $\begin{pmatrix} 0 & i \\ i & 0 \end{pmatrix}$, and $k$ is the label assigned to
    matrix $\begin{pmatrix} i & 0 \\ 0 & -i \end{pmatrix}$; furthermore, we interpret
    $-1, -i, -j, -k$ as the negatives of the corresponding matrices. Under this labeling,
    verify that the group elements satisfy the following rules: The element 1 is the
    identity, and multiplication by $-1$ (in either order) changes the sign of the element.
    Furthermore, $i^2 = j^2 = k^2 = -1$, $ij = k$, $jk = i$, and $ki = j$, while
    $ji = -k$, $kj = -i$, and $ik = -j$. We will henceforth use the notation of this
    exercise when we have occasion to use the group of quaternions.
\end{problem}
\unsolved
